% !TEX program = tectonic
\documentclass[10pt,a4paper,twocolumn]{article}
\usepackage{iftex}
\ifPDFTeX
  \usepackage[T1]{fontenc}
  \usepackage{lmodern}
  \usepackage[utf8]{inputenc}
\else
  \usepackage{fontspec}        % XeTeX/LuaTeX (tectonic) : UTF-8 natif
\fi
\usepackage[french]{babel}
\frenchsetup{StandardLists=true}
\usepackage[a4paper,margin=1.5cm,columnsep=0.6cm,footskip=0.7cm]{geometry}
\usepackage{amsmath,amssymb}
\usepackage{booktabs}
\usepackage{array}
\usepackage{enumitem}
\setlist[itemize]{leftmargin=1.1em,itemsep=0.5pt,topsep=1.5pt,parsep=0pt}
\setlist[enumerate]{leftmargin=1.5em,itemsep=0.5pt,topsep=1.5pt,parsep=0pt}
\usepackage{xcolor}
\usepackage[font=small,skip=2pt]{caption}
\usepackage{titlesec}
\titlespacing{\section}{0pt}{0.75em}{0.3em}
\titleformat{\section}{\normalsize\bfseries}{\thesection.}{0.4em}{}
\usepackage{hyperref}
\hypersetup{colorlinks=true,linkcolor=blue!50!black,urlcolor=blue!50!black}
\setlength{\parindent}{0pt}
\setlength{\parskip}{0.25em}
\raggedbottom
\setlength{\emergencystretch}{1.5em}
\abovedisplayskip=3pt plus 1pt \belowdisplayskip=3pt plus 1pt
\abovedisplayshortskip=2pt \belowdisplayshortskip=2pt
\newcommand{\code}[1]{\texttt{\small #1}}

\title{\vspace{-1.8em}\textbf{Portefeuille des élus du Congrès --- spécification mathématique}\\[2pt]
\large Reconstruction, classement, stratégie de copie et ajustement au marché\\[1pt]
\small Note méthodologique --- \code{05\_Portefeuille\_Membre\_V1.ipynb}}
\author{\small Alice Lemaire --- 6 juillet 2026}
\date{\vspace{-2.6em}}

\begin{document}
\maketitle

\section{Reconstruction du portefeuille d'un membre}
Notations : $i$ = ticker, $t$ = jour de bourse, $P_i(t)$ = cours de clôture.
\begin{itemize}
  \item \textbf{Achat} $\to$ parts, au 1\ier{} jour coté $\ge$ date de transaction :
        $$q = \frac{\text{montant}}{P_i(\tau)},\qquad \tau = \min\{t \ge \code{traded}\}$$
  \item \textbf{Vente} $\to$ retire des parts, \emph{short interdit} :
        $$q^- = \min\!\Big(\tfrac{\text{montant}}{P_i},\; \text{parts détenues}\Big)$$
        le $\min$ règle la vente sans achat connu.
  \item \textbf{Détentions cumulées} (une colonne par ticker) :
        $$n_i(t) = \!\!\sum_{\text{achats}\le t}\! q \;-\!\! \sum_{\text{ventes}\le t}\! q^-$$
  \item \textbf{Valeur du portefeuille} : $\displaystyle V(t) = \sum_i n_i(t)\,P_i(t)$.
\end{itemize}

\section{Rendement d'un membre (\emph{time-weighted})}
On valorise les parts de la \textbf{veille} $\Rightarrow$ un apport n'est jamais compté comme un rendement :
$$r_P(t) = \frac{\sum_i n_i(t-1)\,P_i(t)}{\sum_i n_i(t-1)\,P_i(t-1)} - 1$$
\begin{itemize}
  \item Jour sans position : $r_P(t)=0$ (système ouvert, pas de cash suivi).
  \item Garde-fou : $r_P \leftarrow \operatorname{clip}(r_P,-0{,}5,+0{,}5)$ (winsorisation $\pm 50\,\%$/jour).
  \item Courbes et stats depuis l'origine, avec $\mathrm{NAV}^{\max}(t)=\max_{s\le t}\mathrm{NAV}(s)$ :
        $$\mathrm{NAV}(t)=\!\prod_{s\le t}\!\big(1+r_P(s)\big),\quad
          \sigma_{\text{ann}}=\sigma(r_P)\sqrt{252}$$
        $$\mathrm{Sharpe}=\frac{\overline{r_P}}{\sigma(r_P)}\sqrt{252},\ \
          \mathrm{maxDD}=\min_t\!\Big(\tfrac{\mathrm{NAV}(t)}{\mathrm{NAV}^{\max}(t)}-1\Big)$$
\end{itemize}

\section{Excès, Information Ratio \& ajustement par le marché (\texorpdfstring{$\beta$}{beta})}\label{sec:beta}
Marché = SPY, $r_{\text{SPY}}(t)=\tfrac{\text{SPY}(t)}{\text{SPY}(t-1)}-1$. Série d'excès :
$$x(t)=r_P(t)-r_{\text{SPY}}(t),\qquad
  \mathrm{IR}=\frac{\overline{x}}{\sigma(x)}\sqrt{252}$$
Test de significativité (unilatéral, « bat-il le SPY ? ») :
$$t=\frac{\overline{x}}{\sigma(x)/\sqrt{n}}=\mathrm{IR}\sqrt{\tfrac{n_{\text{jours}}}{252}}\;\;>\;1{,}645$$

\smallskip
\textbf{Ajustement par le marché.} L'IR soustrait \emph{tout} le marché $\Rightarrow$ il suppose $\beta=1$. Un portefeuille agressif
($\beta>1$) affiche alors un excès positif en marché haussier \emph{sans talent}. On régresse plutôt
les rendements quotidiens du membre :
$$r_P(t)=\alpha+\beta\,r_{\text{SPY}}(t)+\varepsilon(t)$$
\textbf{Pente} $\beta$ (exposition) et \textbf{ordonnée} $\alpha$ (talent), par moindres carrés :
$$\beta=\frac{\operatorname{Cov}(r_P,\,r_{\text{SPY}})}{\operatorname{Var}(r_{\text{SPY}})},\qquad
  \alpha=\overline{r_P}-\beta\,\overline{r_{\text{SPY}}}$$
$$\operatorname{Cov}=\overline{(r_P-\overline{r_P})(r_{\text{SPY}}-\overline{r_{\text{SPY}}})},\;\;
  \operatorname{Var}=\overline{(r_{\text{SPY}}-\overline{r_{\text{SPY}}})^2}$$
$\beta$ = pente du nuage de points quotidiens $(r_{\text{SPY}},r_P)$ ; $\alpha$ = son rendement un jour où
le marché est plat.

\textbf{Ratio d'appréciation} = l'IR \emph{du résidu} $\varepsilon=r_P-(\alpha+\beta\,r_{\text{SPY}})$
(le talent risque-ajusté, marché retiré) :
$$\text{appraisal}=\frac{\alpha}{\sigma(\varepsilon)}\sqrt{252}$$
\begin{itemize}
  \item $\beta=1$ : appraisal $=$ IR (identiques) ;
  \item $\beta\neq1$ : l'IR mélange exposition et talent, l'appraisal \textbf{isole} le talent ;
  \item classer par appraisal (au lieu de l'IR) retire le \textbf{levier de style}.
\end{itemize}

\section{Éligibilité \& tests multiples}
\begin{itemize}
  \item \textbf{Éligibilité} : $n_{\text{trades}}\ge 10$ et $n_{\text{jours}}\ge 126$ ($\sim$6 mois).
  \item \textbf{Tests multiples} : sur $M$ membres testés à $5\,\%$, on attend $0{,}05\,M$ « significatifs »
        par hasard. Correction de Bonferroni : seuil par test $\tfrac{0{,}05}{M}$, soit
        $t \gtrsim 3{,}5$ pour $M\simeq 223$.
\end{itemize}

\section{Classement walk-forward (point-in-time)}
À chaque coupe de fin d'année $Y$, on ne garde que l'information $\le 31/12/Y$ :
$$\mathrm{classement}(Y)=\big\{\text{éligibles} \;\wedge\; t>1{,}645\big\}\ \text{triés par } \mathrm{IR},$$
calculés sur $r_P\big|_{\le Y}$. La sélection $\mathrm{topK}(Y)$ prend les $K$ premiers (repli si $<K$).
Aucun regard vers le futur : $n_i(t)$ ne dépend que des trades $\le t$.

\section{Stratégie top-\texorpdfstring{$K$}{K} : la bonne combinaison}
On sélectionne fin $Y$ et on suit le panier en $Y{+}1$. Poids \textbf{initiaux} :
$$w_k=\frac1K \qquad\text{ou}\qquad w_k=\frac{\mathrm{IR}_k}{\sum_j \mathrm{IR}_j},\qquad \sum_k w_k=1.$$

\textbf{Écueil --- la moyenne à poids figés est fausse.} Écrire
$r_{\text{strat}}(t)=\sum_k w_k\,r_{m_k}(t)$ avec $w_k$ \emph{constants} revient à \textbf{rebalancer
chaque jour} vers $w_k$ ; or « copier et \emph{tenir} » laisse les poids \textbf{dériver} avec la performance.

\textbf{Combinaison correcte (buy-and-hold).} On compose la NAV de chaque membre depuis le début de l'année,
puis on pondère :
$$\mathrm{NAV}_k(t)=\!\!\prod_{\substack{s\le t\\ s\in Y+1}}\!\!\big(1+r_{m_k}(s)\big),\quad
  N(t)=\sum_k w_k\,\mathrm{NAV}_k(t)$$
$$\boxed{\,r_{\text{strat}}(t)=\frac{N(t)}{N(t-1)}-1\,}$$
Les poids \textbf{effectifs} dérivent alors tout seuls :
$$w_k(t)=\frac{w_k\,\mathrm{NAV}_k(t-1)}{\sum_j w_j\,\mathrm{NAV}_j(t-1)}$$
Ils sont remis à $w_k$ à chaque re-sélection annuelle (rebalancement \emph{annuel}, pas quotidien).

\textbf{Exemple ($K=2$, $w=\tfrac12,\tfrac12$).} A = Apple, B = Coca. Jour 1 : Apple $+100\,\%$,
Coca $0\,\%$ ; jour 2 : Apple $0\,\%$, Coca $+20\,\%$.

\begin{center}\small\setlength{\tabcolsep}{5pt}
\begin{tabular}{lrrr}
\toprule
 & A & B & Total\\
\midrule
Jour 0 & 50 & 50 & 100\\
Jour 1 & 100 & 50 & 150 \;$(+50\,\%)$\\
Jour 2 & 100 & 60 & 160 \;$(+6{,}67\,\%)$\\
\bottomrule
\end{tabular}
\end{center}
Au jour 2, la moyenne figée donne $\tfrac12\!\cdot\!0+\tfrac12\!\cdot\!20 = 10\,\%$ (faux) ; la vraie
valeur est $\tfrac{160}{150}-1 = 6{,}67\,\%$, car les poids ont dérivé à
$w_A=\tfrac{100}{150}=\tfrac23,\ w_B=\tfrac13$ :
$$\tfrac23\!\cdot\!0\,\% + \tfrac13\!\cdot\!20\,\% = 6{,}67\,\%.$$

\section{Évaluation annuelle \& ajustement au marché}
\textbf{Excès par année} ($1$ année $=1$ observation, excès composé). Soit
$G_Y(r)=\prod_{t\in Y}\big(1+r(t)\big)$ le facteur composé de l'année :
$$x_{\text{strat},Y}=G_Y(r_{\text{strat}})-G_Y(r_{\text{SPY}})$$
$$\mathrm{IR}_{\text{ann}}=\frac{\overline{x}_{\text{strat}}}{\sigma(x_{\text{strat}})},\qquad
  t=\mathrm{IR}_{\text{ann}}\sqrt{n}$$
avec pour seuil $t_{n-1}\approx 1{,}81$ à $n=11$ années.

\textbf{Ajustement (}$\boldsymbol\beta$\textbf{) de la stratégie.} On applique la même régression
qu'au \S\ref{sec:beta} à la série de la stratégie :
$$\boxed{\,r_{\text{strat}}(t)=\alpha+\beta\,r_{\text{SPY}}(t)+\varepsilon\,}$$
$\alpha$ annualisé $\alpha_{\text{ann}}=(1+\alpha)^{252}-1$ = le talent de la stratégie ; significativité par
$t_\alpha=\alpha/\mathrm{SE}(\alpha)$ (OLS ; un écart-type Newey-West élargirait $\mathrm{SE}$,
l'autocorrélation quotidienne étant réelle). Si $\beta>1$, une part de l'excès brut n'est que de
l'exposition au marché : $\alpha<$ excès brut.

\end{document}
